\documentclass[runningheads]{llncs}

\usepackage[T1]{fontenc}
\usepackage{graphicx}
\usepackage{amsmath}
\usepackage[colorlinks=true,linkcolor=black,citecolor=black,urlcolor=black]{hyperref}

\begin{document}

\title{Characterizing Extended Kalman Filter Fusion of Visual Detection and Odometry Under Sparse Observation: A Multi-Robot Digital Twin Case Study}

\titlerunning{EKF Fusion Under Sparse Observation}

\author{Yan Santos\inst{1} \and Alisson [Sobrenome]\inst{1}}

\authorrunning{Y. Santos and A. [Sobrenome]}

\institute{[Instituição], [Cidade], [País]\\
\email{[email]@[dominio]}}

\maketitle

\begin{abstract}
% TODO: 150-250 palavras. Estrutura: problema (fallback discreto YOLO/odom
% no CERISE) -> método (EKF com R adaptativo por confiança, gating
% Mahalanobis) -> resultado (+23.7% no instante de correção; erro contínuo
% pior que odometria sob cobertura de 0-28%) -> contribuição (validação em
% digital twin ROS2/Gazebo real + caracterização honesta dos limites do
% filtro sob observação esparsa, ancorada em Sinopoli et al. 2004).
\keywords{Extended Kalman Filter \and Sensor Fusion \and Multi-Robot Localization \and Sparse Observations \and Digital Twin}
\end{abstract}

\section{Introduction}
% TODO ~1 página: motivação de information fusion em sistemas multi-robô,
% lacuna específica do CERISE (fallback discreto atual), contribuição.

\section{Related Work}

Our work builds on three layers of prior literature: (i) the classical
formulation and convergence theory of EKF-based localization, (ii) the
theory of Kalman filtering under intermittent or sparse observations, and
(iii) recent multi-robot systems that fuse odometry with visual detection
in realistic simulation or physical testbeds.

\subsection{Classical EKF Localization and Convergence}

Bailey and Durrant-Whyte~\cite{bailey2006slam} provide the standard
tutorial formulation of EKF-based simultaneous localization and mapping,
formalizing the motion model $x_k = f(x_{k-1}, u_k) + w_k$ and observation
model $z_k = h(x_k) + v_k$ that our filter follows directly: odometry
drives prediction, and YOLO-based visual detections drive correction. We
adopt this notation throughout Section~\ref{sec:methodology}.

Dissanayake et al.~\cite{dissanayake2001slam} prove convergence properties
of EKF-SLAM covariance under the assumption of \emph{dense} landmark
observation — i.e., that the filter is corrected regularly and
predictably. This assumption is central to our Discussion
(Section~\ref{sec:discussion}): our overhead-camera setup violates it by
design, since a robot is visually detected in only 0--28\% of odometry
readings depending on scenario, and this violation is precisely what
produces the continuous-error finding we report.

\subsection{Kalman Filtering Under Intermittent Observations}

Sinopoli et al.~\cite{sinopoli2004intermittent} establish that, when
observations arrive as a Bernoulli process, there exists a critical
arrival probability below which the expected estimation error covariance
diverges. Our Scenario~3 (curve/occlusion, 0\% visual correction over the
full recording) sits at the limiting case $p \to 0$ covered by this
theory — the covariance saturation we observe is therefore an expected
consequence of the observation regime, not a symptom of filter
misconfiguration.

Censi~\cite{censi2009geometric} revisits Sinopoli et al.'s result and
shows that the existence and value of this critical threshold depend on
the error metric used to summarize the covariance matrix (Euclidean vs.\
Riemannian manifold of SPD matrices). We report divergence using
$\mathrm{trace}(P)$ (Section~\ref{sec:results}), and note this
metric-dependence explicitly as a scoping remark rather than a claim of a
universal threshold.

\subsection{Recent Multi-Robot Visual-Odometric Fusion}

Khosravi et al.~\cite{khosravi2026cooperative} report a decentralized
cooperative localization scheme for two robots, fusing local EKF
estimates with intermittent mutual LiDAR corrections, validated in Gazebo
and on physical hardware. They report a 34--56\% RMSE reduction from
fusion — the same order of magnitude as our $+23.7\%$ correction-instant
gain — but do not report or discuss a continuous-error regime in which
fusion underperforms the unfused baseline, which is the central,
differentiating finding of our work.

Ahmed and Fr\'emont~\cite{ahmed2026maritime} fuse YOLOv8 detections with a
per-target EKF for decentralized multi-robot obstacle detection and
tracking in a maritime scenario.
% TODO (verificar no corpo do paper, não apenas no abstract, antes de
% publicar): (a) composição exata da equipe multi-robô (UAVs/USV); (b) se
% reportam mAP@0.5 comparável ao nosso 0.995 — se confirmado, citar como
% correspondência externa que apoia nossa leitura de que o achado do erro
% contínuo vem de geometria de sensor, não de detector fraco.

Dobrynin and Garcia~\cite{dobrynin2025turtlebot} present a distributed
Kalman filtering testbed using three TurtleBot3 units on ROS2 — the
closest platform match to our own — but solve relative robot-to-robot
localization without a camera, and therefore do not encounter the
sparse-visual-coverage regime we characterize. We treat this as
complementary rather than competing work, and revisit it as a Future Work
direction in Section~\ref{sec:future-work}.

\subsection{Positioning}

Table~\ref{tab:related-work} summarizes how our work relates to these six
references. Unlike Khosravi et al.~\cite{khosravi2026cooperative}, who
report only a favorable fusion gain, our contribution is the joint
report of both the correction-instant gain and the continuous-error
degradation under sparse coverage, together with a theoretical grounding
(\cite{sinopoli2004intermittent,censi2009geometric}) for why this
trade-off is expected rather than a bug in our implementation.

\begin{table}
\caption{Positioning relative to prior work.}\label{tab:related-work}
\begin{tabular}{|p{2.6cm}|p{2.6cm}|p{2.6cm}|p{2.6cm}|}
\hline
\textbf{Reference} & \textbf{Platform} & \textbf{Reports fusion gain?} & \textbf{Reports sparse-coverage degradation?} \\
\hline
Khosravi et al.~\cite{khosravi2026cooperative} & 2 robots, Gazebo + real & Yes (34--56\%) & No \\
\hline
Ahmed \& Fr\'emont~\cite{ahmed2026maritime} & Multi-robot maritime & Yes (detection-based) & No \\
\hline
Dobrynin \& Garcia~\cite{dobrynin2025turtlebot} & 3 TurtleBot3, ROS2 & N/A (no camera) & N/A \\
\hline
\textbf{This work} & 3 TurtleBot3, Gazebo/ROS2 & Yes ($+23.7\%$) & \textbf{Yes (characterized)} \\
\hline
\end{tabular}
\end{table}

\section{Methodology}
\label{sec:methodology}
% TODO ~2 páginas: estado [px,py,theta], Q de POS_DRIFT_ODOM, R de
% box.conf[0], gating Mahalanobis (com ressalva de Altendorfer & Wirkert
% 2015), COV_CAP como track coasting (Bar-Shalom), pipeline de calibração
% de câmera (Zhang 2000), validação em duas etapas (NEES/NIS sintético +
% bag real).

\section{Experimental Setup}
\label{sec:experimental-setup}
% TODO ~1 página: CERISE, TurtleBot3+Gazebo+ROS2, 3 cenários gravados,
% pacote de reprodutibilidade (bag+README+URDF+params+git SHA), MCAP.

\section{Results}
\label{sec:results}
% TODO ~2-3 páginas: figuras lafusion_trajectory.png,
% lafusion_error_over_time.png, lafusion_nees_nis.png,
% lafusion_covariance_trace.png, lafusion_trajectory_dispersion.png.
% Tabela de erro correction-instant (+23.7%) E erro contínuo (EKF pior
% que odom, com trace(P) saturando em 2xCOV_CAP).

\section{Discussion}
\label{sec:discussion}
% TODO: por que EKF (não UKF/PF/IMM) - ver decisão já tomada; limitação do
% gating por Mahalanobis; framing do achado do erro contínuo como
% confirmação de Dissanayake 2001 + Sinopoli 2004, não como falha.

\section{Future Work}
\label{sec:future-work}
% TODO: (a) LiDAR como 3ª fonte de fusão; (b) incerteza formal de
% projeção/Jacobiano; (c) covariance inflation / Q adaptativo por
% tempo-desde-última-correção; (d) combinar com localização relativa
% distribuída (Dobrynin & Garcia 2025).
%
% PRÓXIMOS PASSOS DE IMPLEMENTAÇÃO (código, não só texto — retomar amanhã,
% ver memória project_lafusion_next_implementation_2026_08_13):
% 1. Métrica alternativa de covariância no gráfico de saturação
%    (lafusion_covariance_trace.png): hoje só trace(P); adicionar det(P)
%    ou maior autovalor como segunda curva, motivado por Censi 2009 (o
%    limiar crítico depende da métrica escolhida) — baixo esforço, reforça
%    rigor metodológico sem tocar no filtro em si.
% 2. NÃO implementar fusão via LiDAR mútuo (Khosravi et al. 2026) nem
%    localização relativa distribuída (Dobrynin & Garcia 2025) antes da
%    submissão — ambos exigiriam retrabalhar o EKF e re-validar NEES/NIS +
%    o resultado já obtido (+23.7%), risco alto para o prazo de 16/08.
%    Ambos continuam como Future Work em texto, não em código.

\section{Conclusion}
% TODO ~0.5 página.

\begin{thebibliography}{8}

\bibitem{bailey2006slam}
Bailey, T., Durrant-Whyte, H.: Simultaneous localization and mapping
(SLAM): Part I. IEEE Robotics \& Automation Magazine \textbf{13}(2),
99--110 (2006)

\bibitem{dissanayake2001slam}
Dissanayake, M.W.M.G., Newman, P., Clark, S., Durrant-Whyte, H.F.,
Csorba, M.: A solution to the simultaneous localization and map building
(SLAM) problem. IEEE Transactions on Robotics and Automation
\textbf{17}(3), 229--241 (2001)

\bibitem{sinopoli2004intermittent}
Sinopoli, B., Schenato, L., Franceschetti, M., Poolla, K., Jordan, M.I.,
Sastry, S.S.: Kalman filtering with intermittent observations. IEEE
Transactions on Automatic Control \textbf{49}(9), 1453--1464 (2004)

\bibitem{censi2009geometric}
Censi, A.: Geometric remarks on Kalman filtering with intermittent
observations. arXiv:0906.1637 (2009)

\bibitem{khosravi2026cooperative}
Khosravi, N., Khosravi, N., Bozorg, M., Bahraini, M.S.: Decentralized
Cooperative Localization for Multi-Robot Systems with Asynchronous Sensor
Fusion. arXiv:2603.12075 (2026)

\bibitem{ahmed2026maritime}
Ahmed, M.F., Fr\'emont, V.: Decentralized Multi-Robot Obstacle Detection
and Tracking in a Maritime Scenario. arXiv:2602.12012 (2026)

\bibitem{dobrynin2025turtlebot}
Dobrynin, D., Garcia, I.T.: A TurtleBot3 Hardware Testbed for Distributed
Kalman Filter Localization. Senior Project, California Polytechnic State
University, San Luis Obispo (2025)

\end{thebibliography}
\end{document}
